% Replace the original appendix with this block.
% Requires the main-paper labels eq:confidence, eq:curriculum,
% eq:margin_reliability, and eq:local_obj, and the updated local-gradient analysis.
% Uses the existing booktabs, array, and graphicx packages.
% Target: approximately three acmart two-column pages; verify in the full project.

\appendix
\section{Implementation and Selection Details}
\label{app:selection-analysis}

\paragraph{Candidate admission.}
The three-round availability schedule is specified in the main paper.
In Round~3, let $\mathcal{G}_i^t$ contain the scores of valid local
easy and medium candidates. A local hard pair is admitted when
\begin{equation}\label{eq:hard_gate_app}
c_i^t(z)\geq Q_{0.5}(\mathcal{G}_i^t).
\end{equation}
The hard gate is inactive if $\mathcal{G}_i^t$ is empty.
After local ranking, public candidates must reach the 25th percentile
of selected local scores; this gate is inactive if no local pair is selected.
Quantiles use the sorted value at zero-based index
$\operatorname{round}((n-1)q)$. Ties use seeded random ordering.
Negative margins are not separately excluded. Deduplication identifies
pairs by their prompt, chosen response, and rejected response; it occurs
after allocation and can reduce the final curriculum size.

\paragraph{Public budget and parameter rationale.}
Equation~\eqref{eq:curriculum} prioritizes local pairs and fills only
the residual budget with eligible public pairs. The public-count cap
is $\lfloor\eta_tN_t\rfloor$, with $\eta_1=0.5$ and
$\eta_2=\eta_3=1$; it is not a bound on the public fraction of the
realized set. The Round~1 cap limits supplementation during cold start.
Later rounds retain local-first allocation and the relative admission
gate without an additional public-count cap. The lower-quartile gate
provides a permissive client-relative screen rather than a fixed
absolute threshold. The symmetric interval $[-0.5,0.5]$ maps zero margin
to $0.5$, and $\epsilon=10^{-4}$ avoids exact endpoint scores.
These are design choices, not empirically optimized thresholds.

\paragraph{Single-field score sensitivity.}
For the mapping $f$ in Eq.~\eqref{eq:confidence}, nonexpansiveness
of clipping gives
\begin{equation}
|f(m+v)-f(m)|\leq\frac{|v|}{U-L}.
\end{equation}
If each margin changes by at most $\nu$, each score changes by at most
$s=\nu/(U-L)$. On a fixed candidate list, a top-$k$ boundary gap
greater than $2s$ preserves membership. A quantile changes by at most
$s$, so a gate decision is preserved when its original threshold
distance exceeds $2s$, provided the reference list is fixed.
Changing $[L,U]$ preserves rankings and quantile decisions when all
compared scores remain in the strictly affine region under both
intervals. Saturation and changing candidate lists can invalidate
this conclusion; these bounds do not establish empirical insensitivity.

\paragraph{Multi-field implementation.}
The released selector also supports pooling available margin fields.
For field $j$, set $p_j=f(m_j)$ and compute
\begin{equation}
c_{\mathrm{pool}}=
\frac{\prod_j p_j}{\prod_j p_j+\prod_j(1-p_j)}.
\end{equation}
The default field list is \texttt{client\_margin},
\texttt{reward\_gap}, \texttt{implicit\_reward\_gap},
\texttt{teacher\_margin}, \texttt{dpo\_margin},
\texttt{margin}, and \texttt{margins}; only finite available values
are used. The scorer writes the same value to
\texttt{client\_margin} and \texttt{dpo\_margin}.
If these are the only fields, pooling gives
$p^2/[p^2+(1-p)^2]$, a strictly increasing function of the
single-field score $p$. Different additional fields can change rankings.
Setting \texttt{BEES\_MARGIN\_KEYS=client\_margin} recovers the
single-field rule. The preceding sensitivity bound applies to $f$;
repeated fields are not independent evidence of preference validity.

\section{Derivations for Margin-Based Selection}
\label{app:preference-validity}

\paragraph{Sufficient condition for margin alignment.}
Use the main-paper definitions of $q_i(z)$, $p_i(z)$, and
$\mathcal{H}_i^t$. For a fixed reference $\rho$, the ideal
KL-regularized reward objective has optimizer
$\pi^*(y\mid x)=\rho(y\mid x)
\exp(a r_i^*(x,y)/\beta)/Z(x)$, assuming a finite normalizer.
This follows by writing the objective as
$\beta\log Z(x)-\beta\operatorname{KL}(\pi\|\pi^*)$.
For $t\geq2$, suppose the current scoring policy $\pi_i^t$ and
retained reference $\rho_i^t$ satisfy, on candidate responses,
\begin{equation}\label{eq:app_ideal_policy}
\log\frac{\pi_i^t(y\mid x)}{\rho_i^t(y\mid x)}
=\frac{a_i^t r_i^*(x,y)}{\beta}
-\log Z_i^t(x)+\xi_i^t(x,y),
\end{equation}
where $a_i^t>0$ and $|\xi_i^t(x,y)|\leq b_i^t$.
Subtracting chosen and rejected terms cancels the normalizer, yielding
\begin{equation}\label{eq:app_margin_approx}
\begin{aligned}
m_i^t(z)&=a_i^tq_i(z)+e_i^t(z),\\
e_i^t(z)&=\beta[\xi_i^t(x,y^w)-\xi_i^t(x,y^l)],
\end{aligned}
\end{equation}
with $|e_i^t(z)|\leq2\beta b_i^t$.
This is a sufficient condition, not a guarantee of finite local
training or aggregation. Round~1 requires a separate alignment
assumption because it uses the initial-policy gap.

\paragraph{Validity and selected-set bounds.}
Under $|e_i^t(z)|\leq\delta_i^t$,
$q_i(z)\geq(m_i^t(z)-\delta_i^t)/a_i^t$.
Sigmoid monotonicity proves Eq.~\eqref{eq:margin_reliability}.
For two candidates,
\begin{equation}
q_i(z_1)-q_i(z_2)
\geq\frac{m_i^t(z_1)-m_i^t(z_2)-2\delta_i^t}{a_i^t},
\end{equation}
so margin separation exceeding $2\delta_i^t$ guarantees their
clean-reward ordering. For the realized nonempty curriculum
$\mathcal{S}=\mathcal{D}_i^t$ with minimum raw margin $\gamma_i^t$,
linearity of conditional expectation gives
\begin{equation}\label{eq:app_selected_errors}
\begin{aligned}
\mathbb{E}\!\left[
\frac{\sum_{z\in\mathcal{S}}(1-C_z)}{|\mathcal{S}|}
\,\middle|\,\mathcal{H}_i^t\right]
&=\frac{\sum_{z\in\mathcal{S}}(1-p_i(z))}{|\mathcal{S}|}\\
&\leq\sigma\!\left(
\frac{\delta_i^t-\gamma_i^t}{a_i^t}\right).
\end{aligned}
\end{equation}
Mutual independence among comparisons is unnecessary.
The budget does not ensure $\gamma_i^t>\delta_i^t$.

\paragraph{Gradient bound.}
For $h_{\vartheta}(z)=\beta\Delta_i^t(z)$, the binary
comparison loss has gradient
$(\sigma(h_{\vartheta}(z))-C)\nabla_{\vartheta}h_{\vartheta}(z)$.
Consequently,
\begin{equation}
\begin{aligned}
&\nabla_{\vartheta}\ell_1
-\mathbb{E}[\nabla_{\vartheta}\ell_{C_z}\mid z,\mathcal{H}_i^t]\\
&\qquad=-(1-p_i(z))\nabla_{\vartheta}h_{\vartheta}(z).
\end{aligned}
\end{equation}
Averaging over $\mathcal{S}$, applying the triangle inequality,
and using $\|\nabla_{\vartheta}\Delta_i^t(z)\|\leq G_i^t$
and Eq.~\eqref{eq:app_selected_errors} proves the main-paper
gradient-discrepancy bound. Both objectives use the same selected
pairs and frozen reference. These conclusions depend on alignment
and do not establish calibration or reasoning correctness;
sequence margins can also depend on response length and formatting.

\paragraph{Conditional local descent.}
View Eq.~\eqref{eq:local_obj} as a function of the trainable
LoRA factors $\phi$. If its gradient is $L_F$-Lipschitz along
a full-gradient step $\phi^+=\phi-\lambda\nabla\mathcal{L}_i^t(\phi)$,
the descent lemma gives
\begin{equation}
\begin{aligned}
\mathcal{L}_i^t(\phi^+)
&\leq\mathcal{L}_i^t(\phi)\\
&\quad-\lambda\left(1-\frac{L_F\lambda}{2}\right)
\|\nabla\mathcal{L}_i^t(\phi)\|^2,
\end{aligned}
\end{equation}
for $0<\lambda<2/L_F$. This concerns a fixed local objective;
it does not cover stochastic optimizer steps or establish
convergence across aggregation and curriculum changes.

\section{Training Configuration}
\label{app:hyperparameters}
Table~\ref{tab:hyperparameters} lists the reported settings.
Mixed precision, gradient checkpointing, and FlashAttention-2
support memory-efficient training.

\begin{table}[t]
\centering
\small
\caption{Training and selection settings. $N_t$ is a target budget;
the realized size may be smaller.}
\label{tab:hyperparameters}
\setlength{\tabcolsep}{3pt}
\renewcommand{\arraystretch}{1.02}
\begin{tabular}{@{}p{0.43\columnwidth}p{0.53\columnwidth}@{}}
\toprule
\textbf{Item} & \textbf{Setting} \\
\midrule
Default backbone & Qwen2.5-7B \\
Hardware & 4 NVIDIA RTX 6000 GPUs \\
Clients & 7 MATH; 6 CodeXGLUE \\
Communication rounds & 3 \\
DPO scale $\beta$ & 0.1 \\
Score implementation & Available-field pooling \\
Clipping $[L,U]$; floor $\epsilon$ & $[-0.5,0.5]$; $10^{-4}$ \\
Budget $N_1=N_2=N_3$ & 1000 pairs per client \\
Public cap $\eta_1;\eta_{2,3}$ & 0.5; 1 (no additional cap) \\
Public / hard quantile & 0.25 / 0.50 (Round~3) \\
Negative-margin exclusion & None \\
Tie-breaking / ordering & Seeded \\
Maximum sequence length & 1024 \\
Learning rate / warm-up & $2\times10^{-6}$ / 0.01 \\
Epochs per round & 1 \\
Global / micro batch size & 8 / 1 \\
Precision / DeepSpeed & bfloat16 / ZeRO Stage~2 \\
LoRA rank / scaling / dropout & 64 / 128 / 0.1 \\
LoRA target modules & \texttt{q\_proj}, \texttt{k\_proj},
\texttt{v\_proj}, \texttt{o\_proj}, \texttt{gate\_proj},
\texttt{up\_proj}, \texttt{down\_proj} \\
\bottomrule
\end{tabular}
\end{table}

\section{Supplementary Experimental Results}
\label{app:additional-results}

\paragraph{Client-level round trajectories.}
Tables~\ref{tab:math_iter_results} and~\ref{tab:codexglue_iter_results}
provide the client-level values underlying the aggregate trajectories.
MATH improves in six of seven subjects from Round~2 to Round~3;
Intermediate Algebra changes from 53.50\% to 53.20\%.
All six CodeXGLUE clients improve over the three rounds.

\begin{table}[t]
\centering
\small
\caption{MATH accuracy (\%) by communication round.
Avg. is the subject macro average.}
\label{tab:math_iter_results}
\setlength{\tabcolsep}{5pt}
\begin{tabular}{@{}lccc@{}}
\toprule
\textbf{Task} & \textbf{R1} & \textbf{R2} & \textbf{R3} \\
\midrule
Algebra & 84.30 & 89.60 & 91.20 \\
Counting & 51.30 & 65.20 & 69.60 \\
Geometry & 50.70 & 59.10 & 60.80 \\
Inter. Alg. & 51.10 & 53.50 & 53.20 \\
Num. Theory & 65.60 & 72.40 & 74.40 \\
Prealgebra & 74.10 & 83.10 & 86.20 \\
Precalc. & 52.20 & 55.30 & 56.40 \\
\midrule
Avg. & 61.33 & 68.31 & 70.26 \\
\bottomrule
\end{tabular}
\end{table}

\begin{table}[t]
\centering
\small
\caption{CodeXGLUE BLEU-4 by communication round.
Avg. is the language macro average.}
\label{tab:codexglue_iter_results}
\setlength{\tabcolsep}{6pt}
\begin{tabular}{@{}lccc@{}}
\toprule
\textbf{Language} & \textbf{R1} & \textbf{R2} & \textbf{R3} \\
\midrule
Go & 8.52 & 9.13 & 12.06 \\
Java & 9.34 & 9.98 & 11.06 \\
JS & 6.32 & 6.49 & 8.53 \\
PHP & 13.85 & 14.23 & 14.68 \\
Py. & 59.47 & 79.80 & 90.81 \\
Ruby & 6.33 & 6.73 & 8.27 \\
\midrule
Avg. & 17.31 & 21.06 & 24.24 \\
\bottomrule
\end{tabular}
\end{table}

\paragraph{CodeXGLUE curriculum visualization.}
\label{app:curriculum-selection}
Figure~\ref{fig:curriculum_code_selection_three_rounds} complements
the MATH visualization with programming-language clients.
The selected distributions vary across clients and rounds,
illustrating personalized selection within the shared availability schedule.

\begin{figure*}[t]
\centering
\includegraphics[width=\textwidth,height=0.34\textheight,keepaspectratio]
{figure/three_round_selected_difficulty_code.pdf}
\caption{Candidate and selected margin distributions on CodeXGLUE,
grouped by client, difficulty, and communication round.}
\Description{Candidate and selected CodeXGLUE margin distributions
across programming-language clients and three rounds.}
\label{fig:curriculum_code_selection_three_rounds}
\end{figure*}
